NRR Comments on the Probability of induction HTML - July 28
(reorganised by example — every example now has its own entry, with
"Suggestions" for what is open and "Recently completed" for what is built.)

# Overall:

my goal is for the illustrative examples to appear as seamlessly and to integrate with the text as seamlessly as possible. Where, for example, number change or colour codes are used, those changes and colours can and should appear directly in the text. controls for those variables should be colour coordinated with that bit of text, when possible.

This should be straightforward when the illustration is just showing one of Peirce's examples. But we will need a different approach when we are using one of my examples to show Peirce's point. I will try to handle that example by example in the comments below.

When it makes sense to do so, sliders are preferred to number entry or drop downs. (obviously, if choosing between 5 non-ordered options, its better to use a list. But for things like ratios or the level of certainty required, etc, sliders are best.

As much as possible, the points should be made visually and understood through interaction. The math and working out should be available too, but is not the point.

Assuming this is being read on a normal computer screen, we should endeavour to have the meat of the interactive visible while still in view of the text that inspires it. (This is another way in which making the example explicitly involve the text is helpful.

to the extent possible, we should strive for continuity, using similar examples and graphics to illustrate similar points. this should help with the rhetorical continuity of the paper. especially the first few sections really do build on one another.

when i put something in [square brackets] i mean for that thing to have a slider or other user input.


++++++++++++++++++++++++++++++++++++++++++++++++++++++++++++++++++++++++++++++++++++++++++++++++++++++++++++++++++++++++++++++++++++++++++++++++++++++++++++

# How to read this file

Each example is one entry, in number order, with the unnumbered new examples
after them. Under each:

  Suggestions        — what is still to do. Your notes are kept in your own
                       words. Anything marked [from the build] is an open item
                       I raised while building, not something you wrote.
  Recently completed — what is implemented and checked in the browser.

State of play. Built: 1, 6, 7, 8, 9, 10, 11, 12, 14, 15, 17, 18, 19, 20, 28,
16, 29, plus new examples 26, 27, 30, 31, 32, 33 and 34.
Plus the earlier pass over all of them — the examples now read as continuations
of the text rather than panels beside it, Peirce's own figures move with the
sliders and revert when an example is shut, and every example carries its number
in the margin.

Still to do: the 1–5 group, 13, 22, 23 and 25. Of those, 13, 22, 23 and 25 are
waiting on your own thinking rather than on me.

Last pass, 4 August: 13, 32, 15, 33, 17, 18, 20 and 16, in that order. Their
notes have moved from Suggestions to Recently completed under each example.


++++++++++++++++++++++++++++++++++++++++++++++++++++++++++++++++++++++++++++++++++++++++++++++++++++++++++++++++++++++++++++++++++++++++++++++++++++++++++++

# By example


## 1

### Suggestions

From the "1 - 5" group. No note written yet beyond what is already built.

### Recently completed

#### the whole thing set as an argument
Rebuilt. The "Peirce's Framework" list, the arrow diagram and the panel of
dropdowns are gone; in their place one sentence, which is the argument itself:

  P(If I were to draw the top card from a well shuffled deck of standard
  playing cards, then that card would be exactly A of any suit)

Yellow for the antecedent, green for the "then", blue for the consequent. Every
choice is a word in the sentence rather than a control beside it — click a word
to change it: top / middle / bottom, well shuffled / brand new, deck of standard
playing cards / piquet pack / trick pack, exactly / higher than / lower than,
and the suit. The rank is a slider, since cycling A to K by clicking would be
thirteen presses.

The two settled cases work as you asked. A brand-new pack determines the card
before the draw — the top of a new pack is the ace of spades, the middle the
king of diamonds, the bottom the ace of hearts, which is the order a new pack
actually comes in — so every card is still shown, being a card in the pack, but
all of them are faded except the one that can come up. A trick pack makes every
one of the fifty-two the ace of spades.

Refined after your look. The words that can be changed are bold with a dotted
rule under them, so they read as clickable, and clicking one opens a short list
under it rather than cycling — with five suits, cycling is four clicks to get
back where you started. The "then" is a stronger green than the washes either
side of it, since it is the part that carries the probability, and the same
three colours are used in the formula: P(A → C) with A in the antecedent's
yellow, the arrow in that green, C in the consequent's blue.

The note under it is now yours: "Probabilities do not apply to events. There is
no probability of C. They apply to arguments, and so always depend on how some
set of antecedent facts relate to the consequent fact."

#### (earlier) — colour coding moved into the text
Peirce's own sentence now carries the three colours: antecedent in ochre,
consequent in blue, consequence in green, on every occurrence including the
plain-prose ones later in the sentence, not just the italicised first use. They
are invisible until the example is opened and fade in with it, on the same rule
as the live numbers — shut, the paragraph is ordinary text.


## 2

### Suggestions

None open.

### Recently completed

#### the rule stated over two arguments
Built to your correction, which was the right one: not one P with two blue
clauses inside it, but two arguments and the rule between them —

  P(A → C1)  1/52  +  P(A → C2)  1/52  =  P(A → [C1 or C2])  2/52

Two sentences, sharing an antecedent (the second reads "the same"), each with
its own consequent set by the same chip-and-menu words as example 1. The
identity underneath carries the live fractions in the same three colours.

The condition is not asserted, it is shown. Set the two consequents so that a
card answers to both — an ace of hearts, and an ace of any suit — and that card
is ringed in the grid and the sum visibly counts it twice: 5/52 where the truth
is 4/52. Peirce's "incompatible consequents" is then something you have broken
rather than something you have been told.


## 3 and 4

### Suggestions

[from the build] Rules 3 and 4 in the same style as 2. The chip-and-menu
machinery is written to be shared (argChip and argMenu are at file scope, not
inside example 1), and 2 is now the pattern to copy. 3 is P(A → C1) x
P(A and C1 → C2) = P(A → [C1 and C2]), which needs the second argument's
antecedent to be the first antecedent with C1 added to it — the one real
difference from 2. 4 is the same with independent consequents, P(A → C1) x
P(A → C2), and wants the same treatment as 2's incompatibility: let the reader
break independence and watch the identity fail.

### Recently completed

Nothing yet.


## 5

### Suggestions

From the "1 - 5" group. No note written yet beyond what is already built.

### Recently completed

#### buttons compacted
The bordered box with two headed rows of long labels is one flat row of short
ones, led by "Peirce's:" and "others:". The rest of 5 is untouched.


## 6

### Suggestions

None open.

### Recently completed

#### two arguments, three grids
Moved. The panel sits under the table of the four combinations rather than
directly under the paragraph, so opening it leaves the statement of the case,
its breakdown, and the demonstration on one screen with the same numbers moving
in each.

Cut back hard on the second pass. No preamble, no metal, no box, no gold and
lead, and the whole "unknown metal" mode with it — the reasoning it carried
about what follows when the two agree is ex7's job, and ex7 does it. What is
left is two sliders and three grids directly beneath them.

The two are A_1 → C and A_2 → C. The grids are right and green, wrong and pink,
with the combined grid in the middle, orange where one is right and the other
wrong. Cells can be sorted into blocks or scattered; the scatter is a fixed
shuffle, so toggling back and forth does not deal them again.

Two tabs. The hundred as expected, and a hundred actual cases — where you can
run them one at a time, and the bulk button then reads "Run the remaining 97"
until the set is finished and offers another hundred. Watching them arrive one
by one is, as you say, the thing that makes the point.

The mismatch you noticed was real. The composed grid counted its own cells while
the decomposed one drew the theoretical figures beside them, and the two grids
had been built from two independent columns — each with the right margin, but
paired arbitrarily — so the joint counts had no reason to agree. There is now
one array of a hundred cases and every grid, count and figure reads from it. The
expected hundred is built from the four combinations rather than from two
margins, so the margins and the joints cannot disagree: at 0.81 and 0.93 it is
75 both right, 6 and 18 mixed, 1 both wrong, and the grids read 81/19 and 93/7
exactly as they should.

One honest residue: 75 of 76 agreements is 0.9868, while the long-run figure is
0.9827. That is what rounding to whole cells costs, and both numbers are on
screen together rather than the difference being papered over.

The paragraph stating the case and the table of the four combinations under it
are wired to the sliders again — that binding was lost when the panel was
rewritten. A quiet "Reset to Peirce's 81 and 93" appears beside the other
buttons once the sliders have left his figures, on the same pattern as 14.

Hovering a cell rings the same cell in all three grids and names the case
beneath them. That is the thing that shows the middle grid is made of the two
beside it rather than being a third grid of its own.

Each of the four combinations in the table now carries its value while the panel
is open — 0.7533, 0.1767, 0.0567, 0.0133, summing to one — set before Peirce's
semicolon and gone again when the panel is shut. The two displayed formulas
further down already had theirs.

While doing that I found a transcription slip in the table: the fourth row was
labelled "and both answer correctly", the same as the first. Its own fraction is
7/100 of 19/100, so it is the both-wrong case; I have made it "and both answer
incorrectly". Easily reverted if you would rather keep the edition as it was.

The convergence chart is back, under the grids and only while cases are being
run: the running proportion right for each argument and, in heavier green, the
proportion right among the cases where they agree, each against a dotted line
at what it is set to be. It keeps growing across sets rather than restarting
with each hundred.


## 7

### Suggestions

None open.

### Recently completed

#### the formula in the text evaluates
Both displayed expressions (the quotient, and the chance 81/19 x 93/7 further
down) gain "= 0.9827" and "= 56.6" while ex7 is open, in green so they read as
the result rather than as part of the expression. Peirce prints them
unevaluated, so shut, the "=" clause disappears entirely.

#### laid out like 8, and wired to it
The two sliders now sit side by side across the top, in the same position they
occupy in 8, with the horizontal bars full width beneath them and a good deal
more room to breathe — the canvas went from 300 to 380 and the chart no longer
shares its row with a column of controls. Headed A_1 → C and A_2 → C like the
rest.

The closing note about building the same quantity cell by cell in the previous
demonstration is cut. It pointed at a mode of 6 that no longer exists.

6, 7 and 8 now drive one another. There is one pair of values held in common and
each panel's sliders are a view onto it, so moving one physically moves the
others; all three were given the same range and step (0.02 to 0.98) so it is
value for value.

The reason it was not working: the mirror only wrote to sliders that already
existed. Set the pair in one panel before the others had ever been opened and
they were built at the defaults, which is exactly the case the feature is for.
Each panel now adopts the common pair as it stands when it is built.

The working under 7's bars is one line, close under them: the bottom margin of
the figure went from three lines to almost none and the box round the working is
gone. It is worth having but it is not what the panel is for.

7's reset became the same quiet "Reset to Peirce's 81 and 93" used elsewhere:
invisible while the sliders are at his figures, and it now resets both panels.


## 8

### Suggestions

None open.

### Recently completed

#### nothing under the sliders but the reading
Everything below the sliders is gone. Neither equation is reproduced: Peirce's
chance equation is printed immediately below the panel and the probability
version immediately above it, and both of those now vary with these sliders, so
the algebra stays where he put it.

Making the printed quotient vary needed a small change to the machinery. That
equation belongs to ex7 as much as to ex8, and a figure in the text could only
have one driver. Drivers can now be registered for a container without that
container being open, and can pull each other along, so the 81 and 93 are
registered once for both panels: whichever is open supplies the figures, and if
both are open the one last touched wins. Opening either one animates both
printed equations and the passage between them.

What is left is the two arguments in A_n→C form with a probability slider and a
chance each, and then the final reading: two number lines, the combination as a
probability and the same combination as a chance, both in green because they are
one quantity and not two.

The chance line has no fixed maximum. It steps up through 1, 2, 5 and their
decades to keep the marker on it — 2, then 5, then 100, then 5,000 as you push
the two probabilities up. Watching the ruler itself give way is Peirce's "any
magnitude, however great" without a word of it being said. Evens is marked while
it is legible and drops off once the scale grows past it.

Checked that both fit on one screen: with the panel open at his figures, the
varying quotient sits at the top of the viewport and the second slider at about
three-quarters down, with the number lines below that.

The two sliders and the two number lines now share one column template and
gutter, so all four tracks have the same left edge, right edge and midpoint.
The lines are inset eight pixels each side of the sliders, which is half a
thumb: it makes the markers travel exactly the thumb's range, so the midpoints
coincide rather than merely the boxes.

The chance ruler is now held at 25 until the chance outgrows it, and steps by
decades above that: 25 up to a chance of 25 : 1, then 50, 100, 500 and so on.
(This needed a small correction to the shared scale helper — its floor was
behaving as a minimum rather than a floor, so a chance of 22 was still jumping
the ruler to 50.)

Kept the odds note — it is a gloss rather than a reproduced equation, and you
had asked for it. Trivial to drop.


## 9

### Suggestions

None open. (The belief slider that was outstanding here has moved to 27.)

### Recently completed

#### two stages
Preamble gone: the opening paragraph, the "quantities involved" box and the
instruction line above the controls. The buttons say what they do.

Stage one, two arguments. The setup from 26 given twice, side by side — a
probability slider, the chance it comes to, and the belief that is its
logarithm — sized to sit in the space the old Argument / Tells table occupied.
Both columns share one chance ruler and one belief ruler, so the two arguments
are read against each other rather than each against itself. The belief line is
the axis the bar chart below is drawn on, so the third row of each column is
literally the bar that argument contributes; moving a probability slider moves
its chance, its belief and its bar together.

Stage two, three arguments or more. Columns will not fit side by side at that
point, so it reformats to a row apiece — the old layout, with the changes from
your note:
 - no text inputs. Each argument is named for what it tells, A_n → C for an
   argument for and A_n → ~C for one against. You were right that ~C is the
   way to write it: the same argument at 40 in 100 for C is 60 in 100 for ~C,
   which is exactly what the rewrite on release already did silently.
 - headings are P(A→C), Ch(A→C), ln Ch(A→C).
 - a half tick under every probability slider, and the slider itself takes the
   colour of the way its argument tells — green above the midpoint, terracotta
   below. Crossing the middle recolours the slider, flips the sign of the
   belief and turns the bar around, all while dragging.
 - the even-chance button now says why nothing happened: an even chance is
   1 : 1, log 1 is 0, it sits at the middle of the belief line marked "adds
   nothing", and it can be added any number of times to no effect.

Going back down to two arguments returns to the side-by-side view.

The flip you found is fixed, and by removing rather than mending. An argument
carried a probability and a side, and the two could disagree — so crossing a
half rewrote both on release, which is why the slider jumped and the chance
turned over. An argument now carries one number and no side at all. Which way it
tells is the sign of the logarithm of its chance, and an argument right less
than half the time is simply an argument for ~C. The slider stays exactly where
it is put; the name, the words and the colour change under it as it crosses the
middle, and the chance is stated in the direction the argument actually tells
(0.40 reads A → ~C, tells against, 60 : 40, −0.405). The for/against dropdown is
gone, as you say it should be.

While in here I hoisted the number-line drawing and the growing-ruler scale
into the library — 8, 9 and 26 were about to hold three copies of it.

The even-chance note is trimmed to the clause and the quotation.


## 10

### Suggestions

None open.

### Recently completed

#### re-anchored
The trigger has come off "Suppose we have a large bag of beans" and gone onto
"Suppose the first drawing is white and the next black ... we have great
confidence in this result", stopping before "We now feel pretty sure", and the
panel now opens under that paragraph. No code changed.


## 11

### Suggestions

None open.

### Recently completed

#### the two numbers flattened, and the curve
The two boxes are gone. The numbers now read as one line under the buttons,
first in the blue of the parenthetical it fills in Peirce's sentence above,
second in that one's red, with the count and the range on the line below. The
graph takes the full width. The closing note under the table is cut.

A tab bar over the graph switches between the drawings and the two numbers as a
curve: a gaussian on the estimate, the probable error shaded as the middle half
of its area, with an arrow across the band. The x-axis is held at 0 to 1 rather
than fitted to the curve, so the narrowing is the thing you see — at ten
drawings the band is a third of the page, at a thousand it is a hairline, and
the first number has hardly moved. The figure sits in the legend rather than
inside the band, since the band closes to nothing; the legend takes whichever
top corner the curve leaves free.


## 12

### Suggestions

None open.

### Recently completed

#### a fix found while dumping the prose
The panel opened saying every bean was worth "-Infinity of belief". Its chance
slider is not in the document until update() has built the block that holds it,
and the first update() reads it before that: missing, the read gives nothing,
and the logarithm of nothing is minus infinity. It corrected itself the moment
anything was touched, which is why it had not been caught. Now defaulted.

#### the arguments as lanes, and 11's two graphs beside them
The block of stacked units is gone. The MBR graphic is now laid out as 9's is:
one lane per argument, black to the right and white to the left, ruled off with
the sum underneath. Two thousand lanes will not fit and would say nothing if
they did, since every bean is worth the same, so a run is three lanes, a
vertical ellipsis carrying "× 1,004", and three more.

The sum bar is on a fixed scale — forty beans across, whatever was drawn — so it
comes out identical in Peirce's two cases, as it must, since balancing reasons
returns the same belief for both. Measured at the pixel it is 69 across at
twenty black and 69 at 1,010 against 990, while the lanes above it go from six
rows to fourteen. An excess past forty is clipped and marked, with the figure
printed underneath either way. (An earlier version scaled it against the
drawings; that made the bar disagree with the very thing it illustrates.)

The right-hand block is now 11's two graphs on a tab, at the same height as the
lanes. The drawings: a real order of drawing — the blacks and whites of the
current setting, shuffled and dealt — so the estimate wanders in and lands
exactly on the observed proportion. The distribution: the same gaussian and
probable error as 11.

Balancing reasons is marked on the drawings chart but not on the distribution,
on your point that it does not express a probability. What the distribution
shows is how an estimate of the bag's proportion is spread, and the rule's
figure is not an estimate of that proportion, so putting it on the axis would
make the two look like rival readings of one quantity. That axis is now labelled
"proportion black in the bag" rather than "P(the hidden bean is black)" for the
same reason. The horizontal line on the drawings chart is open to the same
objection if you want it gone — say so and it goes.

Cut: the "holding the excess fixed" heading, its paragraph and the log-scale
plot under it. The table stays, with your replacement note and the Kasser link.

#### tied into balancing reasons
Buttons across the top, and the two sliders are the draw itself: how many beans,
and what proportion of them came up black. Black and white counts follow from
those, which is how the drawing is described in the paragraph.

The two verdicts sit side by side with a graphic under each. Under balancing
reasons, every bean drawn as a unit, black to the right and white to the left —
first as drawn, and then a second row beneath it with the pairs struck out, so
what the sum actually keeps is a row of its own rather than something to be
picked out of the first. At Peirce's second case that reads "990 against 990
cancel; 20 left over", which is the whole complaint in one picture. Blocks stand
in for beans once there are too many to draw one apiece, and it says so.

Under the proportion, both answers on one probability axis: the proportion the
drawings came to with its probable error as a band, and where balancing reasons
puts the same case, dashed. At 1,010 and 990 that is 0.5050 ± 0.0075 against
> 0.999999, the two of them at opposite ends of the same scale.

On your question about the force each bean contributes — no, there is no way to
calculate it, and I have said so in the panel rather than inventing one.
Balancing reasons needs each drawing to be an argument of some fixed force, and
nothing in the bag or the drawings fixes it, which is exactly why it has to be a
slider. The standing note is gone; the button that sets the force to the observed odds
now sits under its slider, and the explanation appears there when it is pressed
rather than sitting on the page whether or not anyone asked. The point of the
button is that it fails informatively: the force of a single bean turns out to
be a fact about all the other beans, so the arguments are not independent after
all and the adding up is not licensed. That seemed a better answer than a
formula.


## 13 — saturn

### Suggestions

[from the build] 13 was built to land as the refutation of the principle 21
stated. With 21's demonstration gone it now arrives without the positive
statement it was answering. Say the word and I will put a sentence into the new
note to carry that weight, or leave 13 to speak for itself.

### Recently completed

#### coming soon, and back on the page
It was not on the page at all. The trigger and the container came out of the
article when 21 was rewritten — the demonstration 13 answers and 13 itself went
together — so the Saturn paragraph had been sitting there as plain prose with
nothing to open. The trigger is back, on the paragraph exactly as it was
(checked byte for byte against the version before it was removed), and it opens
"Coming soon", the same holding page 22 to 25 use.

The example itself is untouched and still built, above in its own file. It is
one name in a list of five; take "example-ex13" out of that list and the colour
chart returns as it was.


## 14

### Suggestions

None open.

### Recently completed

#### the table as printed, until it is asked to be more
The sets are set the way Peirce sets them: letters on the page, no cell round
each one, and nothing over the groups. The captions are gone — he names what the
groups contain in the paragraph after the table, so a caption was saying it
twice and stopping it from looking like a table in a book.

The two sliders are simply there, side by side so they cost as little height as
possible, with the view tabs beneath them. (They were behind a click for a
round; you were right that it was a door with nothing behind it.)

The chart has lost its frame and its grid. Both were drawn in a warm grey and
together they read as a panel laid over the page — the one thing this table is
not supposed to look like. Axes and bars only now, on white.

The first tab is "List of sets". "Reset to Peirce's figures" has moved out of the
slider row to sit beside the binomial-expansion tab, and still appears only once
the sliders have left his figures.

Both paragraphs around it now follow the sliders. The one before — "The relative
number of white balls in the granary might be anything, say one in three. Then
in one-third of the urns the first ball would be white, and in two-thirds
black..." — is spelled out in words as he spells it, so at 1 : 3 it reads one in
four, one-fourth, three-fourths. The one after was already wired.

The sliders travel. A second pair sits under "In the second group, where there is
one b...", and takes over the moment the first pair goes off the top; only one of
the two is ever on screen, and the second writes straight through to the first,
so they are the same control in effect.

Your worry about the table moving the text around was well founded — at seven
balls drawn the listing runs to thousands of lines, and every step of the slider
was throwing the paragraph thousands of pixels down the page. The paragraph is
now held still: its position is measured before and after each change and the
page scrolled back by the difference. Measured drift is zero at every step from
four balls to seven and back. It has to be an instant scroll rather than the
page's usual smooth one, or the paragraph lurches away and glides back on every
step of the drag.


## 15

### Suggestions

None open.

### Recently completed

Built in the earlier pass over all the examples; no separate write-up of that.

#### the instruction made visible
"Click on any bar of the chart to see the exact probability" is now tinted in
the purple a clicked bar takes, so the cue and what it produces are the same
colour. It goes as soon as it has been obeyed, the panel then having the clicked
outcome to report instead. The style is shared, so 18's identical instruction
reads the same way.


## 16 — hypothesis test

### Suggestions

None open.

### Recently completed

#### the difference measured in probable errors, laid along it
The analysis results box is gone. Every figure in it — both probable errors, the
difference, their sum, the multiple — is in Peirce's own paragraph above, so the
box was the same arithmetic printed a third time.

In its place, under the difference arrow, a second rule perforated into lengths
of e1 + e2: each perforation the two errors end to end, blue then terracotta in
the two groups' own colours, so what the multiple is a multiple OF is a thing
you can see rather than a thing you are told. At the census that is eight whole
perforations and most of a ninth, and they are countable. The reading under it
is "8.7 × (e1 + e2) — every dash is the two probable errors end to end". At his
figures the two errors are 0.00034 and 0.00087, so each dash is about a third
blue and two thirds terracotta, which is itself the answer to why "same sample
size" changes the verdict.

On the confidence level, which is the question you raised. The perforations stay
probable errors whatever the slider says, and the slider goes on widening the
shaded bands. Peirce counts the discrepancy in probable errors and then consults
his table for what such a multiple comes to; if the perforations followed the
level too, the level would be counted twice and the picture would stop agreeing
with the sentence it illustrates — the same double count that was flipping the
verdict before. Above 50% the chart says so in a line under the dashes: "the
bands follow the 90% level; the dashes stay probable errors". Easy to swap if
you would rather have it the other way, and then the multiple in the text has to
move with it.

On your other question — the sum of the probable errors is not a modern
technique. Today two independent errors are combined in quadrature,
root(e1² + e2²), which here is 0.00093 against the sum's 0.00121. Summing them
always overstates the combined error, so it is the conservative move: Peirce's
ten times the sum would be about eleven times the modern quantity. His
conclusion is if anything understated. I have left his arithmetic as his; say
the word and both figures can go on the page together.

#### leads with the graph, and the census works itself out
The graph is first, then the two chart toggles, then the buttons, then the
sliders. The headings and horizontal rules are gone.

Every figure in Peirce's worked paragraph is now live and colour-coded: the two
proportions and their counts in the two groups' colours, p and 2p(1-p) in a
third, and the answers — the difference, "one in a 100", the multiple of the
error, the odds — in a fourth. His whole calculation moves: the fraction whose
square root is taken, the root itself, and both probable errors. Shut, it reads
as he printed it.

The counts are sliders on the exponent, three decades of them, with the count
rounded to three figures. A linear slider over ten million would never land on
the round numbers he quotes; this way "near 1,000,000" and "about 150,000" are
both exactly reachable, and so is the 5,000 the chance case wants.

Two more rows on the printed table above the paragraph, not above the graph:
"probable error, white children — 0.477 root(2 × 0.50 × 0.50 / 1,000,000) =
0.00034" and the same for the colored children, with s and p standing at the
figures in use rather than as letters. p is in the green it wears in the prose,
each count and answer in its own group's colour. They are hidden until 16 is
open, by the same rule 17's custom row uses. The confidence slider is back as it
was, a percentage.

The two chart toggles sit directly under the graph, side by side, out of the
way. The long-run frequency box is gone; that figure is in Peirce's own sentence
already, and his sentence now finishes the job: after "only once out of
[3.0e+13] censuses, in the long run" comes "Meaning the difference is almost
certainly real, and not readily attributable to chance" — likely real though not
as certain in the middle band, readily attributable to chance below it. Tinted,
like 20's reading, so it does not read as his. No bolding.

One thing this turned up on the way. The discrepancy was being measured in
whichever bound the confidence control selected, so raising it flipped the
verdict to "readily attributable to chance". That counts the level twice. The
multiple is now always counted in probable errors, which is the quantity Peirce
counts it in — he gets his ten, and then consults the table for what such a
multiple comes to. The bands follow the confidence level; the multiple and the
verdict do not.

#### the three buttons, and where the live figures depart from his print
(These paragraphs were filed under 12 in the old note; they belong here.)

Three buttons: the 1870 census, same sample size, and a difference due to
chance. "Same sample size" is the interesting one — give the colored count the
same million and the same 0.0105 difference goes from 8.7 times the summed error
to 15.6, because the error was mostly coming from the smaller sample.

Where the live figures depart from his print, they are right and he was
rounding: the root is 1/1,414 rather than 1/1,400, the second probable error
0.00087 rather than 0.0008, and the discrepancy 8.7 times the sum rather than
ten. The odds come out at 3.0e+13 rather than his 10,000,000,000, which is his
table's last row rather than a computed figure — his conclusion holds with a
good deal to spare.


## 17

### Suggestions

None open.

### Recently completed

Built in the earlier pass over all the examples; no separate write-up of that.
Its cascade of samples-with-bounds is the graphic 20 now reuses.

#### the slider, the shading, and Peirce's own four
The confidence slider is a plain percentage from 50 to 99, stepping by one, with
its ends named as 20's are: "50% — the probable error" and "95% — standard
today". It ran in nines before — k of them meaning 100(1 − 10^-k) — so that
every row of his printed table out to the ten-billionth was a position on it.
That put nearly the whole travel of the thumb in country no reader has any use
for and made the levels anyone does want unhittable. The far rows are still
printed above with their bounds worked out; they are simply not places the
slider goes, and clicking one of them no longer moves it, since clamping to 99
would have been a lie about which level was being drawn. The three rows within
range still set it on a click.

Rescale is on by default.

The balls-drawn slider goes down to two and steps by one. It began at ten in
tens, which could not be put at four, and four is the case the passage is about.

The heatmap is in, and it is also the answer to your fifth question. The
vertical lines keep their theoretical heights and take their colour from how
often that result has actually been drawn — pale where nothing has landed,
deepening to a dark violet at the commonest. So the empirical distribution is
shown on the very marks of the distribution it is meant to approach, rather than
in a second chart beside it: the two agree when the darkest lines are also the
tallest, and a short run shows as darkness in the wrong place. Only samples of
the current size are counted, an older record's proportion not sitting on this
lattice. The lines thicken below forty drawn, so that at Peirce's four the five
of them are wide enough to read as a shade.

"Peirce's example: one in three, four drawn" clears the record and sets the
three sliders. It comes with the sentence you asked for, which appears whenever
no sample can hit the truth exactly: "While we can never draw a sample that
reflects the true proportion of 0.333 — four drawings can come out at only five
values, and it is not among them — we can be confident that our error bounds
will nevertheless contain it 39.5% of the time."

The figure is the exact one, summed over the lattice rather than counted off a
run, since it is a claim about the construction. At his own case it comes out
below the level claimed, and the sentence says why rather than glossing it: an
induction from four can land on only five values, the window catches one of them
at this proportion and two at others. That is 20's lattice again, and the same
wording. Where the two rates agree to within a point and a half the clause drops
and the sentence is the plain one.

Checked over three hundred samples at his settings: 120 covered, 0.400 against
the exact 0.395.


## 18

### Suggestions

None open.

### Recently completed

Built in the earlier pass over all the examples; no separate write-up of that.

#### what "tolerably certain" comes to
The explainer at the bottom is cut. It was the enumerated list and the tolerance
printed a second time under the paragraph they came from.

His sentence in the text is the readout now: "Thus we should be tolerably
certain — 92.1 times in 100 — of not being in error by more than one ball in
100." Both the figure and the tolerance move with the sliders.

Which of the two is the handle had to be settled, and the certainty is the
better one: it is the vaguer word, it is the one you asked to have made clearer,
and pinning it lets the tolerance be read off. So the slider that was
"Prediction confidence level: 0.50" now reads "Tolerably certain means: 90 times
in 100", and the tolerance is the smallest whole ball in a hundred the drawings
cover that often.

Counted on the exact binomial, so it is the arithmetic he does in the paragraph
rather than a normal curve laid over it: at one ball in a hundred with a hundred
drawings, no ball, one, or two covers 0.921 of the cases — and his own sentence
comes back word for word, with 92.1 as the figure standing behind "tolerably
certain". That number is his; it is what the six probabilities he prints add to.

"Same sample, an even proportion" is the button. The same hundred drawings and
the same certainty, the proportion moved to a half: it reads "tolerably certain
— 91.1 times in 100 — of not being in error by more than eight balls in 100".
One ball against eight, at the same certainty and the same labour, is the whole
of his remark about extreme values. It also unzooms the chart, so the comb below
is on the whole number line.

The shaded band on the chart is now that tolerance rather than an interval of
its own, so the picture, the sentence and the comb are all about one quantity.

#### the comb, and the probable error of a clicked bar
Under the number line, and only on the whole number line — zoomed to five
standard errors round an extreme proportion the teeth are all of a length and
say nothing, so the band is not reserved at all and the panel says where the
comb is instead. Unzoomed, the probable error at every proportion: a comb whose teeth are 0.477 root(2p(1-p)/s) to scale, longest in
the middle and vanishing at the two ends, with the tooth at the current
proportion picked out in the colour its line takes above. The caption gives both
figures — at his own settings, "the probable error at each proportion — ±0.0067
here, ±0.0337 at a half". Five times as much room at the middle for the same
number of drawings, drawn rather than asserted. At a half it says instead that
this is as wide as it ever gets.

On your question about the clicked bar: yes, and it is the relevant one, but not
in the way the bar might suggest. Clicking now draws a band across the top of
the stem — the estimate that sample gives, plus and minus the probable error
worked at the proportion the sample came out at, which is the only proportion an
inquirer holding it would have. The panel says whether that band covers the real
proportion. It also says what the band is not: the height of the bar is how
often that sample turns up, and the band is an error on the estimate; they are
different quantities on different axes, which is exactly the confusion the
question is pointing at.

## 19

### Suggestions

None open.

### Recently completed

#### relaid out to match 20
Slider and buttons share one row, so the hundred Cretans spread the full width
in four rows as they do in 20. Then the text, then the chart: boxes on top,
reading in the middle, graph below. The preamble and the running tally are now
one block in that middle position rather than a paragraph above and a panel to
the side.

Colours are shared with 20 — liars muted red, truth-tellers muted green — so the
hundred Cretans read alike across the two.


## 20 — cretan induction

### Suggestions

[from the build] The sample is capped at twenty rather than forty. Past that the
hundred is too small a population for the bound to hold.

### Recently completed

#### what "approximated" comes to, and the other half of the worst case
"Seventeen or eighteen instances" and "±0.0818" were answering to different
things — the first read as a number of samples, the second worked from the
sample size. Instances are the names: his ±1/6 is 0.477 root(2 × ¼ / 5) =
0.1508 = 1/6.6, which needs five examinations, and he lists five names in the
sentence before. Five samples of one would give the same figure; the formula
knows only how many Cretans have been looked at, not how they were grouped. So
nothing was rewired, and what was missing was said instead.

His sentence now gives the closeness in the units the example counts in:
"probably approximated to within fifteen in a hundred by an induction from five
or six instances". At seventeen it is eight in a hundred. That is what "a decent
approximation" means here, as a number rather than as a word.

His adverb is left alone. "Probably" is a modal about how often the bound holds,
which the confidence slider already governs three clauses later; the closeness is
a different quantity, so it is added rather than swapped for. Say the word and it
becomes an adverb ladder — roughly, fairly, closely — instead.

The figure is his worst case throughout, since that is what the sentence is
about and, until the reveal, the only bound there is. Which shows its own point
when the confidence slider moves: five instances at 95 per cent is "within
forty-four in a hundred", i.e. nothing at all.

And the other half of "even in the worst case", which he leaves implicit and the
page did not have: away from a half the same bound is bought with fewer
instances. Setting 0.477 root(2q(1-q)/n) equal to the worst case at s gives
n = 4sq(1-q), so after the reveal the report adds "At that proportion two
instances would buy the bound that five do at a half". Checked at 0.14 and at
0.88: both come to two.

#### the names are the sample
His list is the sample size, so the two are now one control. The sentence he
prints — "Minos, Sarpedon, Rhadamanthus, Deucalion, and Epimenides, are all the
Cretans I can think of; But these were all atrocious liars, Therefore, pretty
much all Cretans must have been liars" — has four live spans in it: the names,
the clause saying how they were got, the clause saying how many lied, and the
conclusion.

A hundred ancient Greek names, women and men through the roll, his own five at
the front in his own order. Who lies is settled behind the boxes and nowhere
else, so a name takes its colour — the muted red or green the boxes wear — only
once its box has been uncovered. The sentence and the hundred boxes are then the
same fact read twice and cannot disagree. Each box carries its own name on
hover, so a name in the sentence can be found in the roll by hand.

The mixing slider is what takes the sentence away from him. Unmixed, the draw is
the front of the roll, the front of the roll is his five, and they are all
liars — so at the setting his own paragraph describes, his sentence comes out
word for word, and it is only the mixing that makes it something else: "Kalliste,
Menandros, Sosthenes, Timandra, and Eurymedon, are five Cretans drawn at random;
But three of these five were atrocious liars, Therefore, about three-fifths of
the Cretans must have been liars". Before any sample is taken the clause is his
own again, the list then being the front of the roll, which is what the clause
says it is. Shut, the whole sentence is his.

The conclusion follows the draw rather than standing as his caricature over a
sample that does not support it. All liars gives his "pretty much all"; anything
else gives the proportion drawn, spelled the way he spells fractions.

"Take five samples" is in, between one and a hundred.

The slider now reads "Cretans named in each sample", since that is what it is:
how many times one draws is what the buttons do.

On the proportion changing who lies — there is no proportion slider here any
more. It is hidden and drawn afresh by "New Cretans", which is the point of the
rebuild: the enquirer does not know it. Say the word if you would rather have it
settable, but the boxes would then be showing you the answer.

On the reveal numbers: checked again after the names went in. They come from the
same state the boxes do — one induction reported with its own bound, the average
reported separately with no ± claimed — and the two rates the sentence gives,
the worst-case one and the one fitted at the real proportion, still agree with
the exact hypergeometric figure printed beside them.

#### the hundred Cretans
Rebuilt to the note. A hundred boxes reading "?", liars and truth-tellers behind
them in a hidden proportion, as with the bag in 10. A sample uncovers that many
boxes and gives its estimate and bound. Samples accumulate. "Reveal the
proportion" turns the boxes over and colours the record. The old
running-proportion plot is gone; the samples themselves now show the
convergence.

Two plots side by side. The cascade from 17 on these Cretans: every sample with
its bound round it, grey while the proportion is hidden, green and red after,
and the count in the title. Beside it the distribution of estimates, carrying
our best estimate and its probable error as an error bar across the top, so a
biased draw shows as a pile in the wrong place.

Nothing sits in a box beside the example any more. Peirce's sentence now runs
on: after "not be in error by more than [1/6.6]" comes ", meaning an error bound
of ±0.1508", and after his full stop a sentence reporting the run. Before the
reveal it gives the induction and whether the draw earned its bound; after it,
the real proportion, how far out we were, and how often the bound held. It is
set a shade off the body ink so it does not read as Peirce's own, and each
figure wears the colour of the slider that moves it — sample size red, bound and
level green, the proportion the gold its line takes in both charts, the estimate
the purple of its marker.

Why the caught rate came out near a third rather than a half. Two things, and
only one of them was a fault. The fault: the sentence read "our estimate was
0.402 ± 0.1508" where 0.402 was the average of five thousand samples, whose own
probable error is nearer ±0.003. That pinned a one-sample bound to a pooled
number, so "out by 0.002, inside the bound" looked far too easy beside a rate of
34.5 in 100. It now reports one induction — the last, the one whose five are
ringed in the boxes — with its bound, and gives the average separately with no ±
claimed. The histogram's marker follows suit: it is the last induction and its
band, so you can watch that band cover the gold line or miss it.

Not a fault: the rate itself. Two things move it off the claim, and the sentence
now names both, with the numbers, because the first was missing and it is
usually the larger.

The bound is worked out at p = 1/2, where the spread is widest. At any other
proportion the window is wider than that proportion warrants, so it catches more
often than claimed. Your seventeen-instance case: the worst-case bound is
±0.0818, but at the real proportion of 0.90 the probable error is only ±0.0491.
A window 1.7 times too wide is caught 0.796 of the time rather than a half. So
yes — the probable error is perfectly calculable, and both figures are now on
the page. The gap between the two rates is the price of not knowing p, which is
exactly why Peirce states the worst case.

What is left over is the lattice. An induction from s instances can land on only
s + 1 values, so the window catches one of them at some proportions and two at
others. At seventeen and 0.90 the fitted bound would still come out at 63 in 100
rather than 50, and that residue is the lattice. At five instances the lattice
dominates and can pull the rate below a half: a window of ±0.1508 round 0.40
catches one of the six values and gives 0.346, round 0.30 it catches two and
gives 0.67. "New Cretans" is worth pressing a few times to watch it swing.

The sentence branches on which of the two is doing the work: if the worst-case
bound is within two per cent of the fitted one it says so and blames the lattice
outright, otherwise it gives both bounds and both rates.

The report is set in a blue tint that wraps across lines, so it reads as
apparatus rather than as Peirce carrying on, and the figures inside it wear
their sliders' colours — sample size red, bound and every rate green, the
proportion gold, the estimate purple — matching his sentence above.

That sentence replaces the readout beside the boxes, the note under the graphs
and the table of bounds, all of which are gone. The boxes now run the full width
in four rows and the plots sit a screen-height from the sliders.

The bound is his worst case, p = 1/2, and while the proportion is hidden that is
the only bound available. The live phrase that swapped the worst-case clause on
a truth slider is therefore gone, and his sentence is printed as he wrote it.
The freed span went to the word "probably", which the confidence slider
rewrites: at 50 it is his own word, at 90 it reads "ninety times out of a
hundred". Moving that slider keeps the samples; the bands widen and the misses
turn green.

The confidence slider reads "50% — the probable error" at one end and "95% —
standard today" at 95.

On "all I know": one slider rather than a toggle plus an ordering control. The
roll is laid out liars-first, since the Cretans anyone remembers are the
memorable ones, so the two controls collapse into one. Its ends name themselves
— "the five that come to mind" and "a random five", tracking the sample size.
Unmixed, the same five come up every sample, the estimate never moves and none
of the bands catch anything. Mixed, the same arithmetic works.

At five instances the counted rate does not sit on a half, and should not: the
estimate can land on only six values, the bound is set for the worst case, and
drawing five distinct Cretans from a hundred is narrower than five draws from
Crete. The readout gives the exact rate for a fair draw of that size from these
hundred (hypergeometric), so the counted figure has something to be checked
against. At five it runs about 0.35 to 0.67 depending on where the truth falls
between the six values; by twenty it is near a half; at 90 per cent it comes out
about 0.94 against 0.936 exact.

Colours: liars muted red, truth-tellers muted green, uncovered boxes pale. Same
pair applied to 19, so the hundred Cretans read alike across the two.


## 21

### Suggestions

None open.

### Recently completed

#### the example gives way to your text
The indifference demonstration is gone and your Sciential / Ignorantical /
Laplacian paragraph stands in its place, behind a blue rule as a remark of yours
rather than a sentence of Peirce's. It opens on a click of "except that this
introduces the thoroughly unclear idea of cases equally possible", the same
sentence that used to open the demonstration, so the page carries only Peirce
until it is asked to carry more. No collapsing header on this one — there is
nothing to collapse, and calling it an interactive example would promise
something it does not do. It keeps 21 in the margin. Quotation marks and the
citation on the Laplacian Probability passage are set as typography rather than
as typed; nothing else of your wording is touched.

Worth knowing: 13, the Saturn colour chart, was built to land as the refutation
of the principle 21 stated. With 21 gone it now arrives without the positive
statement it was answering. Say the word and I will put a sentence into the new
note to carry that weight, or leave 13 to speak for itself.


## 22

### Suggestions

Should be combined with 23. see below.

### Recently completed

Nothing yet.


## 23

### Suggestions

link it to this text:
"Now, if there be any way of enumerating the possibilities of Nature so as to make them equally probable, it is clearly one which should make one arrangement or combination of the elements of Nature as probable as another, that is, a distribution like that we have supposed, and it, therefore, appears that the assumption that any such thing can be done, leads simply to the conclusion that reasoning from past to future experience is absolutely worthless. In fact, the moment that you assume that the chances in favor of that of which we are totally ignorant are even, the problem about the tides does not differ, in any arithmetical particular, from the case in which a penny (known to be equally likely to come up heads and tails) should turn up heads m times successively. In short, it would be to assume that Nature is a pure chaos, or chance combination of independent elements, in which reasoning from one fact to another would be impossible; and since, as we shall hereafter see, there is no judgment of pure observation without reasoning, it would be to suppose all human cognition illusory and no real knowledge possible. It would be to suppose that if we have found the order of Nature more or less regular in the past, this has been by a pure run of luck which we may expect is now at an end. Now, it may be we have no scintilla of proof to the contrary, but reason is unnecessary in reference to that belief which is of all the most settled, which nobody doubts or can doubt, and which he who should deny would stultify himself in so doing."

i have forgotten the idea for this, so will hold off for now.

### Recently completed

Nothing yet.


## 24

### Suggestions

No note written yet.

### Recently completed

Nothing yet.


## 25

### Suggestions

No note written yet — listed as still to do, waiting on your own thinking.

### Recently completed

Nothing yet.


## 26 (new) — the logarithm of the chance

### Suggestions

None open. (The felt-intensity slider that was outstanding here has moved to
27.)

### Recently completed

#### the logarithm of the chance
Anchored where you asked, on "Now, there is one quantity which, more simply than
any other, fulfills these conditions; it is the logarithm of the chance." Three
number lines of the same length on three scales — probability (p/n), chance
(p/(n−p)), log(chance) — in one grid, so a point on one reads straight down to
the others. Top line is the slider; the other two markers track it with no lag.

The paragraph above is now a readout of that slider. Peirce walks the whole
scale in it, and each clause owns a stretch of the travel, so dragging lights
the sentence describing where you are and fades the one you have left:

  p ≥ 0.985     absolute certainty, an infinite chance, an infinite belief
  0.90 – 0.985  when there is a very great chance, belief ought to be intense
  0.56 – 0.90   as the chance diminishes the feeling of believing should diminish
  0.44 – 0.56   until an even chance is reached, where it should completely vanish
  0.015 – 0.44  when the chance becomes less, a contrary belief springs up
  p < 0.015     as the chance almost vanishes, contrary belief tends to infinity

The ranges live in the markup beside the prose they belong to, not in the
script, so a sentence and its stretch of the slider cannot drift apart. Shutting
the example clears the paragraph back to plain text. Checked that the whole
paragraph and the three lines sit on one screen together. The highlight is a
26% wash rather than the 10% used elsewhere, so it carries across a paragraph
of justified text.

The chance line runs to 100 rather than 10, ticked every ten and labelled at 0,
50 and 100 — a tick per unit would be a smear at that width. Evens is no longer
marked on it, since at 1 in 100 it sits on the origin; it is still marked at the
centre of the log line below, which is rather that line's point. Off-the-line
now begins above a probability of about 0.99 instead of 0.91.

The table of worked values under the sliders is cut.

The value column is a fixed width now rather than max-content, so the readings
changing length ("1 : 1" to "999 : 1 — off the line") no longer resize the
column and shunt the tracks sideways mid-drag. Measured stable at 287px across
the range.


## 27 (new) — which way the connection runs

### Suggestions

None open.

### Recently completed

#### which way the connection runs
Anchored on "The kernel of it is that the conjoint probability of all the
arguments in our possession ... must be intimately connected with the just
degree of our belief in that fact". Your paragraph opens it, as written (I took
out a doubled "of" in "the nature of of this intimate connection").

Two argument sliders give a proper intensity of belief, shown on a belief line
running −6 to +6 with an even chance at the middle. Under it, on the same scale,
a second line that is yours to move, and four dispositions to park it at —
believe just what you ought, stubborn, sceptical, over enthusiastic.

The demonstration is that nothing happens. Move the lower slider through its
whole range and the proper intensity above does not stir; the readout says so in
as many words once you have touched it. Each line also carries the probability
its belief corresponds to, so the gap is legible as a difference in probability
and not only as a difference in nats. Verified that the proper value is
character-for-character identical before and after the belief is dragged.

"Update your belief" is set as text rather than as a button — it is the
sentence the three completions finish — and "to its proper state" is selected
from the start, so something moves on the first drag:

  to its proper state   snaps to whatever the arguments warrant
  stubbornly            moves the right way, at 66% of the evidence's speed
  enthusiastically      moves the right way, at 133% of it

The rate applies to how far the evidence has just moved, not to the gap standing
open. Applied to the gap, all three end up in the same place however slow the
rate, because a drag fires a hundred times and a hundred part-steps arrive where
one whole one does. Applied to the change, the dispositions separate properly:
take the arguments from Peirce's pair down to two even chances and the proper
intensity falls 4.037 to nothing, while the stubborn believer is still at +1.37
and the enthusiast has overshot to −1.33.

The belief is held as a real number and only written to the slider, never read
back off it — rounding each step to the slider's own resolution biased the
accumulation enough, at a stubborn two thirds, to move the belief further than
the evidence had. The slider itself is finer now, 0.01.

Neither intensity shows a probability. Both are unit-less, as your own paragraph
says, and printing one against either was quietly suggesting it could be read
back out.

Labels are Proper Intensity of Belief and Your Intensity of Belief.

The closing line is now the instruction: move the probability sliders to see how
they impact the proper intensity of belief; move the Your Belief slider to see
how much it impacts any worldly probability.

Earlier fixes:
 - The marker on the proper line was invisible. It was being drawn and
   positioned correctly, but its track had no colour key, and the marker had no
   colour of its own to fall back on, so it painted nothing. The marker now
   carries a default colour so this cannot happen silently again, and the
   proper line is keyed green like the other derived readings.
 - The dispositions were in fact wired — they were setting the slider all
   along — but with no visible marker to move against, and with stubborn and
   sceptical moving it only a little, they read as inert. The pressed one is
   now marked until you move the slider by hand.

This closes the outstanding items on both 9 and 26 — it was the same one-way
street both times, and it belongs here rather than in either of them.

#### the note marked, the instruction moved
Your added paragraph is set behind the same blue rule. "Move the probability
sliders to see how they impact the proper intensity of belief. Move the Your
Belief slider to see how much it impacts any worldly probability." has come out
of the readout at the bottom and now sits on its own line at the end of the
preamble, where it can be read before the sliders rather than after them. The
readout keeps the two intensities and the gap between them.


## 28

### Suggestions

None open.

### Recently completed

#### one urn for each proportion (new)
Anchored on "One urn contains all white balls, another one black and the rest
white ... until an urn is reached containing only black balls."

Nothing is calculated. One slider for the number of balls, and the whole row is
there: n+1 urns, all white at one end, all black at the other, one for every
proportion between, each labelled with its ratio. The passage describes an
arrangement and the arrangement is worth seeing before anything is said about
drawing from it. Drawn in the text rather than in a canvas, so it scales with
the page.

#### which urn was this paper drawn from?
Peirce's urns are still there, on a radio button. The second view is the one
from your note: 101 urns, the kth holding letters of which k in a hundred are
consonants, and the paper rewritten in v and c only — every letter in the
article outside this example, in the vowel and consonant colours. Reading is
shared with 29, so the same head serves both.

A flat prior over the 101, then the posterior as bars, with the paper's own
proportion marked. At a thousand letters the weight has gathered on urn 62
against the paper's 61.9. The closing note is the point: nobody drew this paper
from an urn, the urns were supposed at the outset with each as likely as the
next, and the answer is only ever as good as that supposition.

The urns are a second chart directly under the first, sharing its x limits and
its left and right margins, so bar k of one stands over bar k of the other: the
weight above, the urn it belongs to below. Each of the 101 is divided at its own
proportion, consonant blue from the bottom and vowel gold above, so the boundary
runs as a straight diagonal from all vowel at one end to all consonant at the
other. Both ends are labelled in place rather than by a legend, the x-axis is on
the lower chart only, and the paper's own proportion is the same red line
through both.

#### 28 and 29 — giving the article back
Fixed. Two holes: the inner header collapses the content without touching the
container's class, so nothing was watching it, and scrolling away was not
handled at all. Both are now covered by one test — open, not collapsed, and (for
v/c) still on screen.

Both modes now drop on scroll-away, as you asked: the v/c replacement and the
shading alike give the article back the moment their example leaves the screen,
and take it again when it comes back.

Also from that note: the posterior in 28 carries a dashed line at its mean,
running down through the urn chart below it, and the urn it lands on is ringed.
The legend reads "the urns say urn 62" and the text with it: "the weight has a
mean at 61.9, so the urns say this paper came from urn 62 — ringed below". Since
the two charts share an x-axis you can look straight down from the peak to what
that urn holds.

#### balls per urn capped at six
Both "balls in each urn" sliders now stop at six: 30's, where the count of urns
is (r+1)^n and the arithmetic was what slowed down, and 28's alongside it, since
the two rows are meant to be read against each other and a cap on one only would
have put them out of step.


## 29

### Suggestions

None open.

### Recently completed

#### watching the tide (new)
Anchored on "the solution involves the conceptualistic principle that there is
an even chance of a totally unknown event."

It opens at m = 0, which is the whole objection: you have just arrived at the
bay, you have seen the tide do nothing at all, and the rule already answers a
half. The panel says so in as many words. Watch the tide rise and (m+1)/(m+2)
creeps up the line towards one — ten rises puts it at 0.9167 — and "Arrive at
the bay" puts you back at the beginning, where it answers a half again having
seen nothing. The paragraph spelling that out is cut — the marker sitting on the
"even chance" tick before anything has been observed says it.

The half was not sitting over its own tick. Two-line tick labels take the width
of their longer line, and without a text-align the short line sat at its left
edge. Fixed for every number line on the page, not only this one.

#### is the next letter a consonant?
The tide is gone. The rule of succession is now asked about the next letter of
this paper, which has a definite proportion of consonants that nobody knows
before looking. At nothing read it answers a half, and says so: the paper is in
fact 0.619 consonant, so the answer is not merely unfounded but wrong. Read and
the same rule walks over to the truth.

Buttons read one, ten, a hundred, or up to the end of whatever paragraph you are
looking at; "read nothing yet" puts it back to a half, having unlearned
everything. The bar shows the counted proportion as a fill, the rule's answer
marked above it in green, the paper's own proportion below in red. Under that,
the reading itself: the opening as Peirce set it, then an ellipsis, then the
last stretch, vowels gold and consonants blue.

While it is open the article shows its own letters in those colours — faintly
where the reading has not got to, strongly where it has, so the reading head is
a visible boundary running down the page. There are 28,181 letters in the prose
paragraphs, so they are wrapped a block at a time as blocks come into view
rather than all at once. Headings, the apparatus inside examples, the granary
table and the footnotes are not indexed, and neither is anything inside a live
span, since a driven value would change the letter count and put every later
index out.

#### giving the article back
See the entry of that name under 28 — it covers 28 and 29 together.


## 30 (new) — filling the urns from the granary

### Suggestions

None open.

### Recently completed

#### filling the urns from the granary
Anchored on "Suppose we had an immense granary filled with black and white balls
well mixed up; and suppose each urn were filled by taking a fixed number of
balls from this granary quite at random."

Two sliders. The granary's composition, shown as a bar rather than as balls —
there is no end of them, so there is nothing to count. And the number of urns,
which is shared with 28: 28 counts balls and has one more urn than balls, 30
counts urns and puts one fewer ball in each, so the two rows always hold the
same urns and can be read against each other. Move either slider and the other
example follows.

Put one ball in each urn at a time, or fill them, then sort them whitest first.
Defaults are four balls and five urns.

The second slider counts balls, not urns, and the urns follow from it: sharing
the granary out exactly over n balls takes (r+1)^n urns to hold every way they
can fall, once each. At one white in three that is 9, 27, 81 for two, three and
four balls. Which is better than a free count of urns — every setting is now an
exact share-out rather than an arbitrary subset of one.

Past four balls there are more urns than can be set out, so the row is replaced
by the arithmetic: "Sharing the granary out exactly over 5 balls an urn would
take 3^5 = 243 urns — more than can be set out here."

The totals sit under the row: 1, 8, 24, 32, 16 at four balls, which is the
table's own group counts, and 1, 6, 12, 8 at three.

Sorting now matches the table's own order — wwww, wwwb, wwwb, wwbw, wwbw, wbww,
wbww, bwww, bwww — so the row and the listing below it read the same way. (My
first attempt sorted with the whites last, since "b" comes before "w" in the
alphabet; the key is 0 for white and 1 for black.)
The paragraph after this one — "say one in three ... one-third ... two-thirds" —
follows the granary slider, spelled the way Peirce spells it. Those figures used
to be wired to the table further down; they belong here, where the granary is.

The sorting is the point. 28 is the arrangement as stipulated, one urn of every
proportion exactly once. Sort these and you get nothing so tidy: a heap near the
granary's own proportion and few at the ends, which is the distribution the
table below actually counts. The closing paragraph saying so is cut — the row
says it.

#### balls per urn capped at six
Both "balls in each urn" sliders now stop at six: 30's, where the count of urns
is (r+1)^n and the arithmetic was what slowed down, and 28's alongside it, since
the two rows are meant to be read against each other and a cap on one only would
have put them out of step.


## 31 (new) — what the first two balls tell you

### Suggestions

None open.

### Recently completed

#### what the first two balls tell you
Anchored on "Now suppose two balls were drawn from one of these urns and were
found to be both white, what would be the probability of the next one being
white?"

The same eighty-one urns. Draw two at random, or set them — white-white,
white-black, black-white, black-black. The two drawn are shown as large circles
beside a third, dashed and marked with a question mark: what is being asked
about. Choosing a pair keeps only the urns that agree and dims the rest, with
the third ball of each survivor ringed.

Then those third balls are gathered out of the urns and set in a row of their
own, whites first, with the count under them. The four groups are 9, 18, 18 and
36 urns — utterly different sizes — and the row is a third white in every one of
them. A small chart under it puts the four side by side at their true lengths,
each split at the same place.
Cut after your look: the four-group bar chart. The gathered row of third balls
was already carrying that, and the chart was saying it a second time. Peirce's
own conclusion — "Hence the colors of the balls already drawn have no influence
on the probability of any other being white or black" — is not printed in the
panel either. It lights up where he wrote it, in the paragraph the panel hangs
from, on the same mechanism example 1 uses for antecedent and consequent.


## 32 (new) — which arguments are worth balancing

### Suggestions

[from the build] Where it hangs. It is provisionally on a note of its own after
12, set behind a blue rule as an editorial remark — the same treatment 21 and 27
are waiting for. Easy to move once you have picked a home for it.

[from the build] The placeholder on 32's first tab still says "example to come";
it can be pointed at 34 whenever you want it wired up.

### Recently completed

#### the weighted view held back
The tab bar is out, so the panel opens on the rule as stated and there is no
second view to find. Nothing behind it is unpicked — the shrinkage, its column
in the table and the whole of its reading are still there and still working;
the two lines of markup that make the tabs are commented out beside a note
saying so, and putting them back returns it.

One sentence had to move with it. The reading under the rule ended by sending
the reader to the weighted tab, which would now be a door with nothing behind
it. It ends instead on the point that tab was there to sharpen: the rule did not
leave anything out, the line you drew did.

#### which arguments are worth balancing
The constructive half of 12. Several arguments bear on one conclusion, none of
them handed to us: each is established by testing it, so each arrives with a
probable error as well as a probability, and the rule has nowhere to put the
second number.

How well an argument is established is laid on a ladder — three tests at one
end, a thousand at the other — rather than sampled, because a random spread now
and then hands out a set that all clears the line at once and there is nothing
to see. Which rung an argument sits on is shuffled against how strongly it
tells, so being well established and telling strongly stay independent; that is
what stops the threshold from working as a proxy for strength. The probabilities
are not laid out at all, they come out of the testing, which runs in front of
you over about a second when the button is hit.

Left: each argument as what the testing left, a centre and a spread, drawn as a
ridgeline on one P axis. Admitted ones in the colour of the way they tell, ones
too vague in grey and tagged "left out". Right: the balance itself, laid out as
9 lays it, on the admitted arguments only, with the excluded lanes named but
empty. Then the table, with P, probable error, chance and log chance for each.

Two things worth knowing. The probable error is Peirce's worst case, from the
number of tests alone, as with the Cretans in 20 — taken at the observed p it
misbehaves badly here, since three tests that all came out the same way give
p-hat = 0 and so an error of nothing at all, which put the worst established
argument on the page down as the best known. And p-hat is given half a trial's
room at each end so an all-one-way run has a chance rather than an infinite one;
that is a continuity correction rather than a prior, and how extreme a chance
the record may claim grows with the testing behind it.

The reading it produces: at five arguments and ±0.080, three are admitted and
the balance puts the conclusion at 0.035, where admitting all five puts it at
0.005. The rule did not leave the other two out. It would have taken them at
full weight; the line you drew left them out, and moving the line moves the
answer.

A second mode on a tab, answering your question about weighting by the probable
error: yes, and here is the shape of it. Each log-chance is pulled towards zero
by tau squared over tau squared plus sigma squared — the share of the spread
among the arguments that is real rather than their own noise. A rule tested a
thousand times keeps a weight of 1.00; one tested thirteen times gets 0.75; one
tested three times gets 0.41. The full bar is drawn in outline behind the shaved
one, so what the weighting took is a gap you can see rather than a number to be
trusted. On a typical run it moves the balance from 0.976 at face value to
0.969.

Zero is the right place to shrink towards, and not by accident: zero log-chance
is an even chance, which is the one thing Peirce says can do nothing towards
reenforcing others. So a barely tested argument is pulled towards contributing
nothing, by degrees rather than by the cliff the threshold gives.

One correction since first writing this. Tau squared has to be the second
moment about zero rather than the variance about the mean, because zero is where
the shrinking goes. Taken about the mean it got the most ordinary case badly
wrong: five strong arguments that agree with one another have almost no variance
between them, so tau squared came out at nothing and five well-established rules
were weighted away to an even chance. About zero it behaves — at high testing
effort every weight now sits at 0.98 to 1.00, which is what it should be.

It is labelled just for fun, in your words, and the example says why. The rule adds
independent evidence, and you cannot shrink a term in a sum and still claim to
be summing the evidence — weighting is a different operation, one that treats
the arguments as noisy readings of a common quantity to be pooled. That is an
assumption the rule never makes. When it fails outright, which happens when the
whole spread among the arguments is no more than their own noise, tau squared
comes out at nothing, every weight is nothing, and the balance says so on the
chart.

Your replacement text is in on 32's first tab, ending by sending the reader to
the weighted tab. The two now hand off to each other in both directions.


## 33 (new) — when the second number overwhelms the first

### Suggestions

None open.

### Recently completed

#### the button made to work on the probable-error view
It should, and now does. It was not quite dead: it moved the slider, but the
slider steps by 0.002 and one more drawing added to a long run changes the error
by very much less than that, so the step rounded to where it already was and
nothing happened. Over most of the range that is every press.

The error is now held as a real number and only written to the slider, never
read back off it — the same fix 27's belief needed, and for the same reason. The
thumb goes to the nearest step it can show and the reading beside it gives the
exact figure. Every press is one more drawing, exactly: from ±0.010 the count
goes 1,138, 1,139, 1,140, 1,141 while the proportion moves under it. Touch the
slider by hand and it drops the held value; the slider is the handle again.

At the far end the lurch is still the point: at ±0.340, which stands for less
than a single drawing, one press takes it to two drawings and ±0.2385.

With the bag counted exactly there is no drawing left to take, so at an error of
nothing the button goes rather than sitting there doing nothing.

#### when the second number overwhelms the first
Built on your note, between 11 and 12, on the "it is true that when our
knowledge is very precise" paragraph, which had no example on it. Two modes on a
tab, since you wanted both.

"Set the probable error" is your spec. The second number is the quantity under
discussion, so it is the one in your hand, and what follows from it — how much
drawing it would take to have earned an error that size — is printed beside it.
±0.010 is 1,138 drawings; ±0.200 is three. At nothing the bag is counted exactly
and the curve has no width at all. Past ±0.337 the error is wider than a single
drawing could ever leave it, which is to say nothing has been drawn, and that is
where the passage lands. "Take one more sample" steps the slider down by exactly
one drawing's worth: from ±0.340 it goes to ±0.238, a visible lurch, where at
±0.010 it does not move at all.

"Set the drawings" is the same claim on the bag of 10 and 11, where the error is
not chosen but earned. The two disagree about nothing; they differ in which
number is the handle.

Peirce's sentences light as you pass through the regimes he names: very precise
below ±0.02, very slight above ±0.12, and no knowledge at all past ±0.337.
Between them nothing lights, since he names no case there and the two numbers
are simply of a size. The highlighting goes out when the example is shut.

The end of the range is the point. There is no curve, because there is no fact.
The flat line is still drawn — but as the conceptualist's answer and labelled as
his — with "entirely indefinite" underneath, so the disagreement is staged
rather than asserted. The note beside it: a half here is not a measurement that
came out at a half, it is the absence of any measurement written in the same
notation, and the notation is what makes it look like a fact.

Two details worth knowing. The regime is decided on the unrounded number of
drawings — an error of ±0.44 stands for six-tenths of a drawing, and rounding
first called that one drawing and never reached the case the passage is about.
And the curve's spread is the one that goes with the probable error being used
(pe is 0.6745 of it exactly) rather than one read off the observed proportion,
which collapses to nothing the moment a short run comes out all one way and
leaves no curve to draw at just the point where there is least to go on.

Cut on your instruction: the opening paragraph, and the note that appeared in
the "very slight" regime. Peirce's own lit sentence carries that one.


## 34 (new) — a proportion of what?

### Suggestions

None open.

### Recently completed

#### a proportion of what?
The sequel to 27, opening from "cannot be used to determine any probability" in
27's own text, with its panel underneath 27's.

Peirce's two rules drawn as counts. A1 has 100 cases with 81 shaded; A2 has 100
with 93. Each probability is the shaded part of its own grid, countable. The
balance of the two is 0.9827, which is not the shaded part of either. The only
grid it could be the shaded part of is a third one, the cases where both
antecedents held, and nobody has counted it.

You can then go and count it, and the two arguments barely constrain what you
find. At an overlap of twenty the proportion can be anything from 0.650 to
1.000, every value consistent with 81 in 100 and 93 in 100. The rule's 0.9827 is
one of them, near the top. Nothing in the arguments picked it out; independence
did.

Two ends worth the visit. At an overlap of none the two antecedents never occur
together, the third grid is empty, and there is nothing to take a proportion of
at all. At an overlap of 88 — the largest the counts allow — the answer is
pinned at exactly 81/88 = 0.9205, and 0.9827 is not it. Past 88 the shared cases
would have to hold more consequents than A1 contains, so the two arguments
cannot be about more than 88 of the same cases: they provably do not share a
population.

The pooled figure is dealt with in passing: 174 in 200, or 0.87, is countable
but moves when you observe A2 more often, so it is a proportion of how much you
looked rather than of what you found.

This is the example 32 wanted to link to. Same fact seen twice — no common
population, so nothing to pool there and nothing to be a proportion of here.


++++++++++++++++++++++++++++++++++++++++++++++++++++++++++++++++++++++++++++++++++++++++++++++++++++++++++++++++++++++++++++++++++++++++++++++++++++++++++++

# Unnumbered examples


## "can be used to determine any probability"

### Suggestions

This one needs some work... for now let's hide it.

### Recently completed

Nothing yet.


## Blackberries

### Suggestions

"The relative probability of this or that arrangement of Nature is something which we should have a right to talk about if universes were as plenty as blackberries, if we could put a quantity of them in a bag, shake them well up, draw out a sample, and examine them to see what proportion of them had one arrangement and what proportion another. "

not sure yet.

### Recently completed

Nothing yet.


## New example — Socrates, and the two kinds of reasoning

### Suggestions

"All our reasonings are of two kinds: 1. Explicative, analytic, or deductive; 2. Amplifiative, synthetic, or (loosely speaking) inductive. In explicative reasoning, certain facts are first laid down in the premises. These facts are, in every case, an inexhaustible multitude, but they may often be summed up in one simple proposition by means of some regularity which runs through them all. Thus, take the proposition that Socrates was a man; this implies (to go no further) that during every fraction of a second of his whole life (or, if you please, during the greater part of them) he was a man. He did not at one instant appear as a tree and at another as a dog; he did not flow into water, or appear in two places at once; you could not put your finger through him as if he were an optical image, etc. Now, the facts being thus laid down, some order among some of them, not particularly made use of for the purpose of stating them, may perhaps be discovered; and this will enable us to throw part or all of them into a new statement, the possibility of which might have escaped attention. Such a statement will be the conclusion of an analytic inference. Of this sort are all mathematical demonstrations. But synthetic reasoning is of another kind. In this case the facts summed up in the conclusion are not among those stated in the premises. They are different facts, as when one sees that the tide rises m times and concludes that it will rise the next time. These are the only inferences which increase our real knowledge, however useful the others may be."

Have a deductive view where the premises are this is socrates and socrates is a man. have a third sort of greyed out premise that is all men are _____

the conclusion can be therefore socrates is _____

have a button that sort of moves a list through the blank spot in the premises, kind of like setting a date on an old iphone or going through a rolodex. just have a list of like 100 things that all men are and have it woosh through them for 2 seconds, and when it lands on one, have it also appear in the conclusion.
the point is that the informatino oin the conclusion is not in the premises.

Show how a problem in probabilities is set up the same way (have a tab under the deduction view to switch it between probabilities and socrates. highlighting this when in probability view "In any problem in probabilities, we have given the relative frequency of certain events, and we perceive that in these facts the relative frequency of another event is given in a hidden way. This being stated makes the solution."

it could be like i am rolling this dice in my dice box
Die like this rolled that way come up six ___ percent of the time.
THerefore, i shoudl expect this one to come up six ___ percent of the time.



then on the induction view, we want to show that we have different facts.

we could maybe use bean bags. premise one would be that we drew these beans  (could have a slider for like 1-10, and as you move it it randomly generates some beans.
and the beans came from that bag.
and the conclusion could be like, the beans in the bag are of the same proportion.

we need a way to show that no fact about any particular bean in your hand has anything to do with a bean in the bag.

### Recently completed

Nothing yet.


## Purple beans

### Suggestions

this text:
"When we draw a deductive or analytic conclusion, our rule of inference is that facts of a certain general character are either invariably or in a certain proportion of cases accompanied by facts of another general character. Then our premise being a fact of the former class, we infer with certainty or with the appropriate degree of probability the existence of a fact of the second class. But the rule for synthetic inference is of a different kind. When we sample a bag of beans we do not in the least assume that the fact of some beans being purple involves the necessity or even the probability of other beans being so. On the contrary, the conceptualistic method of treating probabilities, which really amounts simply to the deductive treatment of them, when rightly carried out leads to the result that a synthetic inference has just an even chance in its favor, or in other words is absolutely worthless. The color of one bean is entirely independent of that of another.

But synthetic inference is founded upon a classification of facts, not according to their characters, but according to the manner of obtaining them. Its rule is, that a number of facts obtained in a given way will in general more or less resemble other facts obtained in the same way; or, experiences whose conditions are the same will have the same general characters.

In the former case, we know that premises precisely similar in form to those of the given ones will yield true conclusions, just once in a calculable number of times. In the latter case, we only know that premises obtained under circumstances similar to the given ones (though perhaps themselves very different) will yield true conclusions, at least once in a calculable number of times. 25We may express this by saying that in the case of analytic inference we know the probability of our conclusion (if the premises are true), but in the case of synthetic inferences we only know the degree of trustworthiness of our proceeding. As all knowledge comes from synthetic inference, we must equally infer that all human certainty consists merely in our knowing that the processes by which our knowledge has been derived are such as must generally have led to true conclusions."


Modes for analytic and synthetic. synthetic will then split into conceptualist and materialist.

#### Analytic

[slider] % of beans in this bag are purple
These [n] beans were drawn [randomly / selectively] from that bag
:  In the long run, [n] / total beans will approach [the correct percentage].

start the slider at 100% and have it real "all" (when it goes to 0 have it read none). and have the conclusion just follow, these beans will be purple.  as the slider goes into not certainty, modify the conclusion to what is above. note there is no probability about this example, there are however many purple beans there are in the sample. (we could do like the black bean hidden thing above, but have it hide the sample. if the beans are draw selectively, have a note saying that this changes the probability, its like the ace of spades deck abouve, or a magician drawing your card.

if your sample size is the total number of beans, then that becomes also a kind of certainty.

#### Synthetic Conceptualist

needs to make this point: "On the contrary, the conceptualistic method of treating probabilities, which really amounts simply to the deductive treatment of them, when rightly carried out leads to the result that a synthetic inference has just an even chance in its favor, or in other words is absolutely worthless. The color of one bean is entirely independent of that of another."

not sure how to best do this...

#### Synthetic (Materialist)

(no note yet)

### Recently completed

Nothing yet.


## New example — the tides, and the treatises

### Suggestions

"Most treatises on probability contain a very different doctrine. They state, for example, that if one of the ancient denizens of the shores of the Mediterranean, who had never heard of tides, had gone to the bay of Biscay, and had there seen the tide rise, say m times, he could know that there was a probability equal to (m + 1) / (m + 2) that it would rise the next time. In a well-known work by Quetelet, much stress is laid on this, and it is made the foundation of a theory of inductive reasoning."

could have bayes theorem pop up and show how you need values for the porbability of events?

### Recently completed

Nothing yet. (29 already carries the rule of succession itself, on the letters
of this paper rather than the tide.)
